\documentclass[UTF8,12pt,a4paper,fontset=none]{ctexart}
\usepackage[a4paper,left=1.82cm,right=1.82cm,top=1.72cm,bottom=1.82cm,headheight=15pt]{geometry}
\usepackage{fontspec,xeCJK}
\setmainfont{Liberation Serif}
\setsansfont{DejaVu Sans}
\setmonofont{DejaVu Sans Mono}
\setCJKmainfont[BoldFont=Noto Serif CJK SC Bold]{Noto Serif CJK SC}
\setCJKsansfont[BoldFont=Noto Sans CJK SC Bold]{Noto Sans CJK SC}
\setCJKmonofont{Noto Sans Mono CJK SC}
\usepackage{amsmath,amssymb,mathtools,bm}
\usepackage{booktabs,longtable,tabularx,array,multirow}
\usepackage{enumitem,xcolor}
\usepackage[most]{tcolorbox}
\usepackage{fancyhdr,titlesec,hyperref,cleveref,lastpage}
\usepackage{tikz,float,caption,listings}
\usetikzlibrary{arrows.meta,positioning,calc,fit,backgrounds,shapes.geometric}
\definecolor{mainblue}{RGB}{39,82,138}
\definecolor{deepblue}{RGB}{25,55,95}
\definecolor{softblue}{RGB}{235,243,252}
\definecolor{softgreen}{RGB}{238,248,241}
\definecolor{softorange}{RGB}{255,246,232}
\definecolor{softgray}{RGB}{245,246,248}
\definecolor{warnred}{RGB}{160,48,48}
\hypersetup{unicode=true,colorlinks=true,linkcolor=deepblue,urlcolor=mainblue,citecolor=deepblue}
\setlength{\parindent}{2em}
\setlength{\parskip}{0.36em}
\setlength{\tabcolsep}{5pt}
\renewcommand{\arraystretch}{1.2}
\allowdisplaybreaks[4]
\emergencystretch=3em
\sloppy
\pagestyle{fancy}\fancyhf{}\fancyfoot[C]{\small 第 \thepage\ 页，共 \pageref{LastPage} 页}\renewcommand{\headrulewidth}{0.4pt}
\titleformat{\section}{\Large\bfseries\color{deepblue}}{\thesection}{0.65em}{}
\titleformat{\subsection}{\large\bfseries\color{mainblue}}{\thesubsection}{0.55em}{}
\titleformat{\subsubsection}{\normalsize\bfseries}{\thesubsubsection}{0.5em}{}
\numberwithin{equation}{section}
\newcommand{\R}{\mathbb{R}}
\newcommand{\E}{\mathbb{E}}
\newcommand{\Prob}{\mathbb{P}}
\newcommand{\Loss}{\mathcal{L}}
\newcommand{\norm}[1]{\left\lVert #1\right\rVert}
\newcommand{\ip}[2]{\left\langle #1,#2\right\rangle}
\newcommand{\tr}{\operatorname{Tr}}
\newcommand{\diag}{\operatorname{diag}}
\newcommand{\rank}{\operatorname{rank}}
\newcommand{\softmax}{\operatorname{softmax}}
\newcommand{\argmin}{\operatorname*{arg\,min}}
\newcommand{\argmax}{\operatorname*{arg\,max}}
\newcommand{\sg}{\operatorname{sg}}
\newcommand{\KL}{D_{\mathrm{KL}}}
\newcommand{\dd}{\mathrm{d}}
\newtcolorbox{keybox}[1][]{enhanced,breakable,colback=softblue,colframe=mainblue,boxrule=0.8pt,arc=1.5mm,left=2mm,right=2mm,top=1.2mm,bottom=1.2mm,title={#1},fonttitle=\bfseries}
\newtcolorbox{examplebox}[1][]{enhanced,breakable,colback=softgreen,colframe=green!45!black,boxrule=0.7pt,arc=1.5mm,left=2mm,right=2mm,top=1.2mm,bottom=1.2mm,title={#1},fonttitle=\bfseries}
\newtcolorbox{notebox}[1][]{enhanced,breakable,colback=softorange,colframe=orange!70!black,boxrule=0.7pt,arc=1.5mm,left=2mm,right=2mm,top=1.2mm,bottom=1.2mm,title={#1},fonttitle=\bfseries}
\newtcolorbox{criticalbox}[1][]{enhanced,breakable,colback=red!3,colframe=warnred,boxrule=0.8pt,arc=1.5mm,left=2mm,right=2mm,top=1.2mm,bottom=1.2mm,title={#1},fonttitle=\bfseries}
\lstset{basicstyle=\ttfamily\small,breaklines=true,frame=single,backgroundcolor=\color{softgray},numbers=left,numberstyle=\tiny,showstringspaces=false,columns=fullflexible}
\hypersetup{pdftitle={42 ChordEdit 数学过程详解},pdfauthor={OpenAI}}
\fancyhead[L]{\small 42 ChordEdit}
\fancyhead[R]{\small 数学过程详解}
\setlength{\abovedisplayskip}{0.82em}
\setlength{\belowdisplayskip}{0.82em}
\setlength{\abovedisplayshortskip}{0.45em}
\setlength{\belowdisplayshortskip}{0.45em}
\begin{document}
\begin{center}
{\LARGE\bfseries 42\_ChordEdit 数学过程详解}\par
\vspace{0.35em}
{\large 动态最优传输、Chord 控制场与一步稳定性}
\end{center}
\vspace{0.45em}
\begin{keybox}[本文只处理真正需要推导的部分]
本文逐步推导连续性方程、Benamou--Brenier 动能、局部严格凸估计、Chord 公式、核平滑偏差--方差、Jensen 能量收缩、Euler 误差与投影解释。。全文按“公式从哪里来--每个符号是什么--中间步骤为何成立--维度和复杂度如何核对--论文结论依赖哪些假设”的顺序展开；不复述实验表格，也不使用与论文无关的通用背景凑页数。
\end{keybox}
\tableofcontents
\newpage

\section{从条件流差到编辑控制场}
设源提示和目标提示对应的连续流速度分别为
\begin{equation}
 v(x,t,c_{\mathrm{src}}),\qquad v(x,t,c_{\mathrm{tar}}).
\end{equation}
理想的瞬时编辑控制为
\begin{equation}
 \Delta v(x_t,t)
 =v(x_t,t,c_{\mathrm{tar}})-v(x_t,t,c_{\mathrm{src}}).
\end{equation}
但一阶生成器通常只能在合成噪声状态 $z\sim K_t(\cdot\mid x_\tau)$ 上查询。令模型可观测输出为 $Q(z,t,c)$，用时间相关线性映射 $B_t$ 转为统一速度单位，则代理场
\begin{equation}
 \boxed{
 R(x_\tau,t)
 =\E_{z\sim K_t(\cdot\mid x_\tau)}
 \left[B_t\{Q(z,t,c_{\mathrm{tar}})-Q(z,t,c_{\mathrm{src}})\}\right].
 }
\end{equation}
论文把它建模为
\begin{equation}
 R(x_\tau,t)=u_t(x_\tau)+\varepsilon_t,
 \qquad \E\varepsilon_t=0,
\end{equation}
其中 $u_t$ 是真实编辑控制场，$\varepsilon_t$ 是查询噪声与模型误差。

\section{Benamou--Brenier 动态最优传输}
令 $\rho_t(x)$ 表示图像分布沿编辑时间演化的密度。质量守恒要求
\begin{equation}
 \boxed{
 \partial_t\rho_t(x)+\nabla_x\cdot(\rho_t(x)u_t(x))=0.
 }
\end{equation}
其来源可由任意区域 $\Omega$ 的质量守恒推导：
\begin{align}
 \frac{\dd}{\dd t}\int_\Omega\rho_t(x)\dd x
 &=-\int_{\partial\Omega}\rho_tu_t\cdot n\dd S\\
 &=-\int_\Omega\nabla\cdot(\rho_tu_t)\dd x.
\end{align}
因为对任意 $\Omega$ 都成立，得到连续性方程。动态 OT 选择满足边界分布的最低动能路径：
\begin{equation}
 \boxed{
 \min_{\rho,u}\int_0^1\int\frac12\norm{u_t(x)}_2^2\rho_t(x)\dd x\dd t
 \quad\text{s.t. continuity equation}.
 }
\end{equation}
这解释了“低能量”不是经验形容词，而是速度平方关于运输质量的积分。

\section{局部二次估计问题}
直接取 $u_t^{\mathrm{nai}}=R_t$ 会把高频噪声完整带入一步积分。ChordEdit 在窗口 $[t-\delta,t]$ 内把控制近似为常向量 $u$，并最小化
\begin{equation}
 \Phi_t(u;x_\tau)
 =t\norm{u-\widehat u_{t-\delta}(x_\tau)}_2^2
 +\int_{t-\delta}^{t}\norm{u-R(x_\tau,\xi)}_2^2\dd\xi.
\end{equation}
对 $u$ 求梯度：
\begin{align}
 \nabla_u\Phi_t
 &=2t(u-\widehat u_{t-\delta})
 +2\int_{t-\delta}^{t}(u-R_\xi)\dd\xi\\
 &=2(t+\delta)u
 -2t\widehat u_{t-\delta}
 -2\int_{t-\delta}^{t}R_\xi\dd\xi.
\end{align}
令梯度为零，唯一极小点为
\begin{equation}
 \boxed{
 u_t^*(x_\tau)
 =\frac{t}{t+\delta}\widehat u_{t-\delta}(x_\tau)
 +\frac1{t+\delta}\int_{t-\delta}^{t}R(x_\tau,\xi)\dd\xi.
 }
\end{equation}
Hessian 为
\begin{equation}
 \nabla_u^2\Phi_t=2(t+\delta)I\succ0,
\end{equation}
故该解确为严格凸问题的唯一全局最优解。

\section{实用 Chord Control 公式}
采用一阶因果近似
\begin{equation}
 \widehat u_{t-\delta}\approx R(x_\tau,t-\delta),
 \qquad
 \int_{t-\delta}^{t}R(x_\tau,\xi)\dd\xi
 \approx\delta R(x_\tau,t),
\end{equation}
得到
\begin{equation}
 \boxed{
 \widehat u_t(x_\tau)
 =\frac{tR(x_\tau,t-\delta)+\delta R(x_\tau,t)}{t+\delta}.
 }
\end{equation}
权重
\begin{equation}
 \lambda_1=\frac{t}{t+\delta},\qquad
 \lambda_2=\frac{\delta}{t+\delta}
\end{equation}
非负且和为一，所以 $\widehat u_t$ 是两个代理场的凸组合。$\delta=0$ 时退化为 naive field；$\delta$ 增大时平滑增强，但语义响应可能变弱。

\section{核平滑解释}
一般形式可写为单边因果卷积
\begin{equation}
 \widehat u(t)=(K_\delta*R)(t)
 =\int_0^\delta K_\delta(s)R(t-s)\dd s,
\end{equation}
其中
\begin{equation}
 K_\delta(s)\ge0,
 \qquad \int_0^\delta K_\delta(s)\dd s=1.
\end{equation}
若 $R$ 二阶可微，Taylor 展开
\begin{equation}
 R(t-s)=R(t)-sR'(t)+\frac{s^2}{2}R''(t)+O(s^3).
\end{equation}
代回得偏差
\begin{equation}
 \operatorname{Bias}[\widehat u]
 =-\mu_1(K_\delta)R'(t)
 +\frac{\mu_2(K_\delta)}2R''(t)+O(\delta^3),
\end{equation}
其中 $\mu_j=\int s^jK_\delta(s)\dd s$。单边核一般有一阶偏差 $O(\delta)$；若构造一阶矩为零的二阶核，则偏差变为 $O(\delta^2)$，平方偏差为 $O(\delta^4)$。

若观测噪声方差为 $\sigma^2$，独立采样密度与窗口宽度成正比，则平滑方差近似
\begin{equation}
 \operatorname{Var}(\widehat u)=O\!\left(\frac{\sigma^2}{\delta}\right).
\end{equation}
二阶核的均方误差
\begin{equation}
 \operatorname{MSE}(\delta)
 \approx C_1\delta^4+C_2\frac{\sigma^2}{\delta}.
\end{equation}
求导并令零：
\begin{equation}
 4C_1\delta^3-C_2\frac{\sigma^2}{\delta^2}=0,
\end{equation}
故
\begin{equation}
 \boxed{\delta^*\asymp\sigma^{2/5}.}
\end{equation}
这给出窗口宽度的偏差--方差平衡。

\section{$L^2$ 能量收缩的完整证明}
对固定 $t$，把 $S\sim K_\delta$ 看成随机变量，则
\begin{equation}
 \widehat u(t)=\E_S[R(t-S)].
\end{equation}
平方范数是凸函数，由 Jensen 不等式
\begin{equation}
 \norm{\widehat u(t)}_2^2
 \le\E_S\norm{R(t-S)}_2^2.
\end{equation}
对 $t$ 积分并用 Fubini/Tonelli 交换积分次序：
\begin{align}
 \int\norm{\widehat u(t)}^2\dd t
 &\le\int\int K_\delta(s)\norm{R(t-s)}^2\dd s\dd t\\
 &=\int K_\delta(s)\left[\int\norm{R(\tau)}^2\dd\tau\right]\dd s\\
 &=\int\norm{R(\tau)}^2\dd\tau.
\end{align}
因此
\begin{equation}
 \boxed{\norm{K_\delta*R}_{L^2}\le\norm{R}_{L^2}.}
\end{equation}
若核不是 Dirac 脉冲且 $R$ 在核支持内并非常值，平方范数严格凸，故不等号严格。乘以非负密度 $\rho_t(x)$ 并对 $x$ 积分，同样得到 Benamou--Brenier 动能不增。

\section{导数范数与一致性误差}
微分与卷积可交换：
\begin{equation}
 \partial_t\widehat u=K_\delta*(\partial_tR),
 \qquad
 \nabla_x\widehat u=K_\delta*(\nabla_xR).
\end{equation}
由单位质量非负核
\begin{equation}
 \norm{\partial_t\widehat u}_\infty
 \le\norm{\partial_tR}_\infty,
 \qquad
 \norm{\nabla_x\widehat u}_\infty
 \le\norm{\nabla_xR}_\infty.
\end{equation}
显式 Euler 解
\begin{equation}
 \dot x=u(x,t),\qquad
 x_{n+1}^{E}=x_n+h u(x_n,t_n)
\end{equation}
的局部截断误差由轨迹二阶导控制：
\begin{equation}
 \ddot x=\partial_tu+(\nabla_xu)u.
\end{equation}
故
\begin{equation}
 \norm{\tau_{n+1}}
 \le\frac{h^2}{2}
 \left(\norm{\partial_tu}_\infty
 +\norm{\nabla_xu}_\infty\norm{u}_\infty\right).
\end{equation}
定义一致性代理
\begin{equation}
 C(u)=\norm{\partial_tu}_\infty
 +\norm{\nabla_xu}_\infty\norm{u}_\infty,
\end{equation}
则平滑后 $C(\widehat u)\le C(R)$，因此同样的一步步长下局部误差上界更小。

\section{一步稳定性}
若 $u$ 关于 $x$ 的 Lipschitz 常数为 $L=\norm{\nabla_xu}_\infty$，两个 Euler 轨迹误差满足
\begin{equation}
 \norm{e_{n+1}}
 \le(1+hL)\norm{e_n}+\norm{\tau_{n+1}}.
\end{equation}
递推得到全局误差 $O(h)$，常数随 $e^{LT}$ 增长。Chord 平滑不增大 $L$，因此稳定裕度不差于 naive field。这里“可以一步走完”并非 Euler 阶数提高，而是同一一阶离散在更平滑、更低能量场上的误差常数下降。

\section{投影与最低能量解释}
把可实现控制限制在局部 chord 子空间 $\mathcal U_\delta$，令 $P_\delta$ 为 $L^2(\rho)$ 正交投影，则
\begin{equation}
 \widehat u=P_\delta R.
\end{equation}
正交分解
\begin{equation}
 R=P_\delta R+(I-P_\delta)R
\end{equation}
给出 Pythagoras 恒等式
\begin{equation}
 \norm{R}_{L^2(\rho)}^2
 =\norm{\widehat u}_{L^2(\rho)}^2
 +\norm{(I-P_\delta)R}_{L^2(\rho)}^2.
\end{equation}
因此在同一可实现子空间中，投影具有最小平方误差和不高于原场的能量；被去掉的正是与平滑 chord 不相容的高频分量。
% AUGMENTED_DETAILED_MATH_2

\section{连续性方程的弱形式}
对任意光滑测试函数 $\varphi(x,t)$，连续性方程
\begin{equation}
 \partial_t\rho+\nabla\cdot(\rho u)=0
\end{equation}
乘 $\varphi$ 并分部积分得到
\begin{equation}
 \int_0^1\int
 [\partial_t\varphi+\nabla\varphi\cdot u]\rho\dd x\dd t
 +\int\varphi(x,0)\rho_0\dd x
 -\int\varphi(x,1)\rho_1\dd x=0.
\end{equation}
弱形式允许密度不光滑，也是动态 OT 数值离散的基础。

\section{Benamou--Brenier 的动量变量凸化}
令动量
\begin{equation}
 m=\rho u.
\end{equation}
连续性约束变成线性
\begin{equation}
 \partial_t\rho+\nabla\cdot m=0.
\end{equation}
动能
\begin{equation}
 \frac12\int\frac{\norm m^2}{\rho}\dd x\dd t
\end{equation}
对 $(\rho,m)$ 是凸的（perspective of quadratic）。这解释了动态 OT 虽同时优化密度和速度，经过变量替换可形成凸问题。

\section{两点 Chord 公式的能量上界}
实用场
\begin{equation}
 \widehat u=\lambda R_1+(1-\lambda)R_2,
 \qquad 0\le\lambda\le1.
\end{equation}
平方范数恒等式
\begin{equation}
 \norm{\lambda a+(1-\lambda)b}^2
 =\lambda\norm a^2+(1-\lambda)\norm b^2
 -\lambda(1-\lambda)\norm{a-b}^2.
\end{equation}
所以不仅有 Jensen 上界，还得到明确能量下降量
\begin{equation}
 \lambda(1-\lambda)\norm{R_1-R_2}^2.
\end{equation}
两个时间点代理场差异越大，平均带来的高频能量削减越明显。

\section{噪声方差下降}
测量
\begin{equation}
 R_i=u_i+\epsilon_i,
 \qquad \E\epsilon_i=0.
\end{equation}
两点加权估计噪声
\begin{equation}
 \epsilon_c=\lambda\epsilon_1+(1-\lambda)\epsilon_2.
\end{equation}
若独立同方差 $\sigma^2I$，
\begin{equation}
 \operatorname{Var}(\epsilon_c)
 =[\lambda^2+(1-\lambda)^2]\sigma^2I.
\end{equation}
$\lambda=1/2$ 时方差减半。若两次查询噪声相关系数 $\rho$，则
\begin{equation}
 \operatorname{Var}(\epsilon_c)
 =[\lambda^2+(1-\lambda)^2+2\rho\lambda(1-\lambda)]\sigma^2.
\end{equation}
高度相关时平滑的降噪收益变小。

\section{一阶积分近似的误差}
论文用
\begin{equation}
 \int_{t-\delta}^{t}R(\xi)\dd\xi
 \approx\delta R(t).
\end{equation}
Taylor 展开
\begin{equation}
 R(t-s)=R(t)-sR'(t)+\frac{s^2}{2}R''(\xi_s).
\end{equation}
积分得
\begin{equation}
 \int_{t-\delta}^{t}R(\xi)\dd\xi
 =\delta R(t)-\frac{\delta^2}{2}R'(t)+O(\delta^3).
\end{equation}
所以右端点矩形公式误差为 $O(\delta^2)$。若用梯形公式
\begin{equation}
 \frac\delta2[R(t-\delta)+R(t)]
\end{equation}
积分误差可降为 $O(\delta^3)$，但论文的递归先验权重使最终公式不是简单等权梯形。

\section{局部凸目标的 Bayesian 解释}
目标
\begin{equation}
 t\norm{u-u_{prev}}^2+\int\norm{u-R_\xi}^2\dd\xi
\end{equation}
可解释为高斯 MAP。先验
\begin{equation}
 u\sim\mathcal N(u_{prev},(2t)^{-1}I),
\end{equation}
观测
\begin{equation}
 R_\xi\mid u\sim\mathcal N(u,\frac12I)
\end{equation}
的负对数后验正比于该目标。解是按精度加权的后验均值，$t$ 越大越信任历史 chord，$\delta$ 越大越信任新窗口观测总量。

\section{显式 Euler 全局误差常数}
误差递推
\begin{equation}
 e_{n+1}\le(1+hL)e_n+\frac{h^2}{2}M.
\end{equation}
从 $e_0=0$ 迭代 $N=T/h$ 步
\begin{align}
 e_N
 &\le\frac{h^2M}{2}\sum_{j=0}^{N-1}(1+hL)^j\\
 &=\frac{hM}{2L}[(1+hL)^N-1]\\
 &\le\frac{hM}{2L}(e^{LT}-1).
\end{align}
Chord 同时降低 $L=\norm{\nabla u}_\infty$ 和 $M=\norm{\partial_tu+(\nabla u)u}_\infty$，所以减少的不只是步长稳定阈值，也包括全局误差常数。

\section{Proximal refinement 的数学位置}
传输一步得到 $x_{pred}$ 后，论文可选用目标条件模型做
\begin{equation}
 x_{prox}=B_{t_c}Q(x_{pred},t_c,c_{tar}).
\end{equation}
可把它视为近端映射
\begin{equation}
 \operatorname{prox}_{\lambda F}(x)
 =\argmin_y\left[F(y)+\frac1{2\lambda}\norm{y-x}^2\right].
\end{equation}
它在保持靠近 $x_{pred}$ 的同时下降目标语义能量 $F$。该步骤不属于低能量 OT 路径，因此能量评估应在 refinement 前完成，否则会混淆 transport 与语义增强。
\clearpage
\section{动态最优传输的 Benamou--Brenier 形式}
给定源分布 $\rho_0$ 和目标分布 $\rho_1$，二阶 Wasserstein 距离满足
\begin{equation}
 W_2^2(\rho_0,\rho_1)=
 \inf_{\rho_t,u_t}
 \int_0^1\int\rho_t(x)\|u_t(x)\|^2\dd x\dd t,
\end{equation}
约束为连续性方程
\begin{equation}
 \partial_t\rho_t+\nabla\cdot(\rho_tu_t)=0,
 \qquad \rho_{t=0}=\rho_0,\ \rho_{t=1}=\rho_1.
\end{equation}
目标最小化整个概率质量搬运的动能。ChordEdit 的低能量控制思想来自：编辑场不应只追求终点语义，还要避免剧烈、曲折的轨迹。

\section{连续性方程的弱形式}
对光滑测试函数 $\varphi(x,t)$，将连续性方程乘 $\varphi$ 并积分：
\begin{equation}
 \int\varphi\partial_t\rho+
 \int\varphi\nabla\cdot(\rho u)=0.
\end{equation}
分部积分得到
\begin{equation}
 \frac{\dd}{\dd t}\int\varphi(x,t)\rho_t(x)\dd x
 =\int\left(\partial_t\varphi+
 u_t\cdot\nabla\varphi\right)\rho_t\dd x.
\end{equation}
若 $X_t$ 满足 $\dot X_t=u_t(X_t)$，右侧正是
\begin{equation}
 \frac{\dd}{\dd t}\E[\varphi(X_t,t)].
\end{equation}
所以 ODE 粒子流与密度连续性方程描述同一传输。

\section{动量变量使问题凸化}
定义动量
\begin{equation}
 m_t=\rho_tu_t.
\end{equation}
目标变为
\begin{equation}
 \int\frac{\|m_t(x)\|^2}{\rho_t(x)}\dd x\dd t,
\end{equation}
约束
\begin{equation}
 \partial_t\rho_t+\nabla\cdot m_t=0.
\end{equation}
函数 $\|m\|^2/\rho$ 在 $\rho>0$ 上是 convex perspective，因此对 $(\rho,m)$ 联合凸。原变量 $(\rho,u)$ 中的 $\rho\|u\|^2$ 看似双线性，动量替换揭示了凸结构。

\section{局部二次平滑问题的解析解}
设原始编辑场观测为 $R(t)$，希望输出平滑控制 $u(t)$，考虑
\begin{equation}
 \min_u
 \int\|u(t)-R(t)\|^2\dd t+
 \lambda\int\|u'(t)\|^2\dd t.
\end{equation}
Euler--Lagrange 方程
\begin{equation}
 u-R-\lambda u''=0.
\end{equation}
在整条实线上，其 Fourier 变换满足
\begin{equation}
 \widehat u(\omega)=
 \frac1{1+\lambda\omega^2}\widehat R(\omega).
\end{equation}
因此平滑器是低通滤波，高频编辑抖动被衰减。时域核为指数型
\begin{equation}
 K(t)=\frac1{2\sqrt\lambda}e^{-|t|/\sqrt\lambda}.
\end{equation}

\section{因果 Chord 核的归一化}
在线编辑不能使用未来场，只能对过去做因果平均
\begin{equation}
 \widehat u(t)=\int_0^\delta K_\delta(s)R(t-s)\dd s,
\end{equation}
其中
\begin{equation}
 K_\delta(s)\ge0,
 \qquad \int_0^\delta K_\delta(s)\dd s=1.
\end{equation}
若取均匀核
\begin{equation}
 K_\delta(s)=\frac1\delta\mathbf 1_{[0,\delta]}(s),
\end{equation}
则是最近 $\delta$ 时间的弦平均。核越宽，方差越低，但时间滞后越大。

\section{$L^2$ 能量收缩的完整证明}
对归一化非负核，Young 卷积不等式
\begin{equation}
 \|K*R\|_2\le\|K\|_1\|R\|_2.
\end{equation}
因为 $\|K\|_1=1$，
\begin{equation}
 \boxed{\|\widehat u\|_2\le\|R\|_2.}
\end{equation}
也可用 Jensen：
\begin{align}
 \|\widehat u(t)\|^2
 &=\left\|\int K(s)R(t-s)\dd s\right\|^2\\
 &\le\int K(s)\|R(t-s)\|^2\dd s.
\end{align}
再对 $t$ 积分并交换积分次序得到同样结论。平滑场的总动能不超过原始场。

\section{独立噪声下的方差下降}
设
\begin{equation}
 R_k=u+\epsilon_k,
 \qquad \E\epsilon_k=0,
 \quad \operatorname{Var}(\epsilon_k)=\Sigma.
\end{equation}
平均 $m$ 个观测
\begin{equation}
 \widehat u=\frac1m\sum_{k=1}^{m}R_k.
\end{equation}
若噪声独立，
\begin{equation}
 \operatorname{Var}(\widehat u)=\frac1m\Sigma.
\end{equation}
若相邻噪声相关，相关系数 $\rho_{ij}$，
\begin{equation}
 \operatorname{Var}(\widehat u)
 =\frac1{m^2}
 \left(m\Sigma+
 2\sum_{i<j}\operatorname{Cov}(\epsilon_i,\epsilon_j)\right).
\end{equation}
正相关会降低平均的方差收益。

\section{一步 Euler 的稳定性误差}
目标 ODE
\begin{equation}
 \dot x=u(x,t).
\end{equation}
从 $t$ 到 $t-h$ 一步
\begin{equation}
 \widehat x_{t-h}=x_t-h\widehat u(x_t,t).
\end{equation}
真解 Taylor 展开
\begin{equation}
 x_{t-h}=x_t-hu(x_t,t)+
 \frac{h^2}{2}\ddot x(\xi).
\end{equation}
误差
\begin{equation}
 \widehat x_{t-h}-x_{t-h}
 =-h(\widehat u-u)-
 \frac{h^2}{2}\ddot x(\xi).
\end{equation}
平滑控制降低 $\|\widehat u-u\|$ 的随机波动，也降低 $\ddot x=\partial_tu+J_uu$ 的高频项，使大步长更稳定。

\section{两点弦速度与平均速度}
若已知轨迹两端 $x_a,x_b$，弦速度
\begin{equation}
 u_{\rm chord}=\frac{x_b-x_a}{b-a}.
\end{equation}
由基本定理
\begin{equation}
 x_b-x_a=\int_a^bu(x_t,t)\dd t,
\end{equation}
所以
\begin{equation}
 u_{\rm chord}=\frac1{b-a}
 \int_a^bu(x_t,t)\dd t.
\end{equation}
它是沿真实轨迹的平均速度，而非端点处瞬时速度。若轨迹曲率大，弦方向可能偏离局部切向量，但用一步跨越整个区间时，弦恰好对应净位移。

\section{一维 Gaussian 传输例子}
源 $X_0\sim\mathcal N(\mu_0,\sigma_0^2)$，目标 $X_1\sim\mathcal N(\mu_1,\sigma_1^2)$。最优传输映射
\begin{equation}
 T(x)=\mu_1+\frac{\sigma_1}{\sigma_0}(x-\mu_0).
\end{equation}
位移插值
\begin{equation}
 X_t=(1-t)X_0+tT(X_0).
\end{equation}
速度
\begin{equation}
 u_t(X_t)=T(X_0)-X_0
 =(\mu_1-\mu_0)+
 \left(\frac{\sigma_1}{\sigma_0}-1\right)(X_0-\mu_0).
\end{equation}
这条速度在粒子路径上恒定，因此一步 Euler 精确。复杂图像编辑的困难在于真实条件分布间映射不是全局仿射，控制场会随位置和时间变化。

\section{局部二次目标的 Bayesian 解释}
若把原始场观测看成
\begin{equation}
 R=u+\epsilon,
 \qquad \epsilon\sim\mathcal N(0,\sigma^2I),
\end{equation}
并给平滑先验
\begin{equation}
 p(u)\propto
 \exp\left(-\frac\lambda2\int\|u'\|^2\dd t\right),
\end{equation}
则 MAP 估计最小化
\begin{equation}
 \frac1{2\sigma^2}\int\|R-u\|^2\dd t+
 \frac\lambda2\int\|u'\|^2\dd t.
\end{equation}
因此 chord 平滑既可解释为最小能量控制，也可解释为对噪声场做 Bayesian 正则化估计。
\end{document}
